\chapter{Matrices, Eigenvalues and the Tensor Product}\label{ch:linear}

\begin{etymology}
\textbf{Transformation} comes from \emph{transformare}: \emph{trans-},
``across'', and \emph{formare}, ``to shape''. To carry something from one form
to another according to rules.

\textbf{Tensor} comes from the Latin \emph{tendere}, ``to stretch''. William
Rowan Hamilton used the word in 1846 for something else (the length of a
quaternion). Woldemar Voigt gave it its modern sense in 1898, for the quantities
that describe how a crystal stretches and deforms differently in different
directions.
\end{etymology}

Chapter~\ref{ch:hilbert} gave the space where states live. This one is about what
acts on them, and about how the spaces of several qubits combine. Three ideas carry
almost all of quantum computing: linear transformations, which are the gates;
eigenvalues, which are what measurements and many algorithms extract; and the
tensor product, which is where both the power of many qubits and entanglement
come from.

\section{Linear transformations and matrices}

A transformation $T$ between vector spaces is \emph{linear} if it respects their
structure:
\[
  T(\alpha u + \beta v) = \alpha\,T(u) + \beta\,T(v).
\]
Transforming a combination gives the same combination of the transformed
vectors. Most functions are not linear, but the ones that are have a remarkable
property: they are fully determined by what they do to the vectors of a basis.
Write $T(e_1), \dots, T(e_n)$ as columns, and you have a matrix.

\begin{history}{1850--1858 \textperiodcentered{} London}
James Joseph Sylvester coined the word \emph{matrix}, Latin for womb, in 1850,
for the rectangular array from which determinants are born. His friend Arthur
Cayley, a lawyer who did mathematics on the side, published ``A Memoir on the
Theory of Matrices'' in 1858. It treated matrices as objects in their own right,
with addition and multiplication, and defined multiplication so that it composes
transformations.
\end{history}

That is why matrix multiplication has its odd row-times-column formula: it was
designed so that the matrix of ``apply $S$, then $T$'' is the product $TS$. A
quantum circuit is a sequence of gates, so the matrix of a whole circuit is the
product of the matrices of its gates.

The determinant of a matrix measures how it scales areas (or volumes), with a
sign for orientation. A determinant of zero means the transformation flattens the
space and loses a dimension, so it cannot be undone. Unitary matrices always have
determinants of absolute value 1.

\section{Eigenvalues and eigenvectors}

Most transformations mix directions: a vector along $e_1$ may be sent to some
combination of $e_1$ and $e_2$. But some directions are special. They are only
stretched or shrunk, never turned. A vector $v \neq 0$ with
\[
  T v = \lambda v
\]
is an \emph{eigenvector} of $T$, and $\lambda$ is its \emph{eigenvalue}: the
direction stays, the scale changes by $\lambda$.

\begin{history}{1829--1904 \textperiodcentered{} Paris and Göttingen}
Cauchy proved in 1829 that symmetric matrices have real eigenvalues, while
studying quadratic forms. The words came later. In 1904 Hilbert, working on
integral equations, used \emph{Eigenwert}, ``own value'', and called the set of
those values the \emph{spectrum} of the operator. He later said he had chosen
the name with no idea that it would turn out to describe the spectra of atoms.
Twenty years later it did: the energy levels of an atom are the eigenvalues of its
Hamiltonian.
\end{history}

Operators for physical quantities are \emph{Hermitian} ($A = A^\dagger$), and
their spectra are always real, as a measurement result must be. The
\emph{spectral decomposition} says such an operator is entirely given by its
eigenvectors and eigenvalues:
\[
  A = \sum_k \lambda_k \ket{v_k}\bra{v_k} .
\]
In the basis of its eigenvectors, $A$ is diagonal, and diagonal is simple.

Take $A = \begin{pmatrix} 3 & 1 \\ 1 & 2 \end{pmatrix}$. Its eigenvalues are
$(5 \mp \sqrt5)/2$, about $1.382$ and $3.618$. Along the two eigenvectors, $A$
only stretches, by those factors. Every other vector gets turned. Rebuilding
$\lambda_1\ket{v_1}\bra{v_1} + \lambda_2\ket{v_2}\bra{v_2}$ gives back $A$ exactly.

\begin{concept}{Spectrum}
The set of eigenvalues of an operator.

\textbf{What it fixes:} it reduces an operator to the numbers that matter. For a
Hamiltonian, they are the energies; for a measurement, the possible results.

\textbf{What it costs:} for a matrix of size $2^n$, computing the spectrum
classically takes time that grows with $2^n$. Extracting it on a quantum computer
is what phase estimation does, and a large part of useful quantum computing is,
underneath, spectral processing.
\end{concept}

\section{Tensors}

A vector has one index, $v_i$; a matrix has two, $M_{ij}$. A \emph{tensor} can
have any number. But a tensor is not just a table with many indices. It is a table
that changes in a specific way when you change the basis, so that equations
written with tensors look the same in every coordinate system.

\begin{history}{1887--1915 \textperiodcentered{} Padua and Zurich}
Gregorio Ricci-Curbastro developed an ``absolute differential calculus'' in the
1880s and 1890s, a way of writing geometry that does not depend on coordinates.
With his student Tullio Levi-Civita he published the full method in 1900. Almost
nobody used it until Marcel Grossmann showed it to his friend Albert Einstein in
1912. Einstein needed exactly that to describe curved spacetime, and general
relativity, in 1915, is written in tensors.
\end{history}

Quantum computing uses less of this machinery than the name suggests. What it
uses all the time is one operation that builds new spaces out of old ones.

\section{The tensor product}

Given two spaces $V$ and $W$, of dimensions $m$ and $n$, the \emph{tensor product}
$V \otimes W$ is a space of dimension $mn$. If $\{e_i\}$ is a basis of $V$ and
$\{f_j\}$ a basis of $W$, the $mn$ products $e_i \otimes f_j$ are a basis of
$V \otimes W$.

One qubit lives in $\mathbb{C}^2$. Two qubits live in
$\mathbb{C}^2 \otimes \mathbb{C}^2 = \mathbb{C}^4$, with basis
$\ket{00}, \ket{01}, \ket{10}, \ket{11}$. And $n$ qubits live in
$\mathbb{C}^{2^n}$. Dimensions multiply: each new qubit doubles the size of the
space.

That doubling is the source of both the promise and the problem. Describing 50
qubits on a classical computer takes $2^{50} \approx 1.1\times10^{15}$ complex
amplitudes. At 16 bytes each, that is about 18 petabytes of memory, for one state.
This is Feynman's argument from Chapter~\ref{ch:history}, in numbers.

\section{Entanglement}

Some states of two qubits can be written as a product of one state for each:
$\ket{0}\otimes\ket{1} = \ket{01}$, or
$\ket{+}\otimes\ket{+} = \tfrac12(\ket{00} + \ket{01} + \ket{10} + \ket{11})$.
These are \emph{separable}. But most vectors of $V \otimes W$ cannot be split this
way. The simplest example is the Bell state
\[
  \ket{\Phi^+} = \frac{\ket{00} + \ket{11}}{\sqrt 2} .
\]
Try to write it as $(a\ket0 + b\ket1)\otimes(c\ket0 + d\ket1)$: you need $ac$ and
$bd$ to be nonzero and $ad$ and $bc$ to be zero, which is impossible. The two qubits
have no state of their own. Only the pair does.

\begin{history}{1935 \textperiodcentered{} Princeton and Oxford}
Einstein, Boris Podolsky and Nathan Rosen published a paper arguing that states
like this show quantum mechanics is incomplete. Schrödinger answered the same year
and gave the phenomenon its name, \emph{Verschränkung}, which he translated as
\emph{entanglement}. He called it ``not one but rather the characteristic trait of
quantum mechanics''.
\end{history}

\begin{concept}{Entanglement}
A state of several systems that cannot be written as a product of states of each
system. Mathematically, non-separability in the tensor product.

\textbf{What it fixes:} correlations between qubits live in the state itself,
with no extra bits to store them. Every algorithm with an exponential speedup
produces large amounts of it.

\textbf{What it costs:} entangled states are fragile. Disturbing one qubit
disturbs the whole state, and the environment is always disturbing something.
\end{concept}

\section{Measuring entanglement: the Schmidt decomposition}

Any state of two systems can be written as
\[
  \ket{\psi} = \sum_{k=1}^{r} \lambda_k\,\ket{a_k}\otimes\ket{b_k},
\]
with orthonormal $\ket{a_k}$, orthonormal $\ket{b_k}$ and positive coefficients
$\lambda_k$. The number of terms $r$ is the \emph{Schmidt rank}. If $r = 1$, the
state is separable. If $r > 1$, it is entangled, and the more evenly spread the
$\lambda_k$, the more entangled it is.

Finding the decomposition is easy: arrange the four amplitudes of a two-qubit
state as a $2\times2$ matrix and compute its singular values. That is the
\emph{singular value decomposition}, found by Eugenio Beltrami in 1873 and Camille
Jordan in 1874. Erhard Schmidt rediscovered it in 1907 for integral equations,
which is why physicists use his name. For $\ket{01}$ and for
$\ket{+}\otimes\ket{+}$, the singular values are $(1, 0)$: rank 1, separable. For
the Bell state they are $(0.707, 0.707)$: rank 2, as entangled as two qubits can be.

Entanglement, the resource with no classical analogue, turns out to be a question
about the rank of a matrix, a tool from nineteenth-century linear algebra.

\section{Identical particles: bosons and fermions}\label{sec:identical}

The tensor product keeps its factors apart: there is a first qubit and a second
qubit, and $\ket{01}$ is not $\ket{10}$. For two qubits on a chip that is right,
because they sit in different places. For two electrons it is wrong. Every
electron in the universe is identical to every other one. Nothing carries a label
saying ``electron 1'', so swapping two electrons cannot change anything that can be
measured.

Write the swap as an operator, $P(\ket a \otimes \ket b) = \ket b \otimes \ket a$.
For identical particles, $P\ket\psi$ must be the same physical state as $\ket\psi$,
and Chapter~\ref{ch:hilbert} showed that two vectors describe the same state only if
they differ by a global phase: $P\ket\psi = e^{i\theta}\ket\psi$. Swapping twice
puts everything back, $P^2 = I$, so $e^{2i\theta} = 1$, and the phase can only be
$+1$ or $-1$. Nature uses both.
\[
  \text{bosons: } \frac{\ket{ab} + \ket{ba}}{\sqrt2},
  \qquad
  \text{fermions: } \frac{\ket{ab} - \ket{ba}}{\sqrt2}.
\]
Now put both particles in the same state, $a = b$. The boson state becomes
$\sqrt2\,\ket{aa}$, a perfectly good state once normalized. The fermion state
becomes $\ket{aa} - \ket{aa} = 0$, which is not a state at all. Two fermions can
never occupy the same state. That is the \emph{exclusion principle}, and here it
fell out of one minus sign.

\begin{history}{1924--1926 \textperiodcentered{} Dacca, Hamburg, Florence, Cambridge}
In June 1924 Satyendra Nath Bose, a lecturer at the University of Dacca, sent
Einstein a short paper that derived Planck's law by counting photons as truly
indistinguishable. Einstein translated it into German, had it published, and
extended it to atoms. In January 1925 Wolfgang Pauli, in Hamburg, explained the
shells of electrons in atoms with a rule: no two electrons in the same state. In
1926 Enrico Fermi, in Florence, and Paul Dirac, in Cambridge, worked out
independently how particles obeying that rule fill up energy levels. Particles of
the first kind became \emph{bosons}, after Bose, and of the second \emph{fermions},
after Fermi. In 1940 Pauli proved which is which: particles with integer spin are
bosons, particles with half-integer spin are fermions.
\end{history}

The sign matters well beyond atoms, and this book runs into it three times.

\textbf{Matter.} Electrons are fermions, so they cannot all drop into the lowest
level of an atom; they stack up in shells, and chemistry is the result. Photons are
bosons, so they can pile into one state by the billion, and a laser is the result.

\textbf{Superconductors.} Two electrons bound together, a \emph{Cooper pair}, behave
as a boson: swapping two pairs swaps two fermions twice, and $(-1)^2 = +1$. Bosons
can share one state, so below a critical temperature all the pairs of a metal fall
into a single collective state. That is superconductivity, and it is what makes the
qubits of Chapter~\ref{ch:hardware} work.

\textbf{Computation.} The minus sign is where the cost of simulating electrons
hides. A classical simulation of many electrons has to keep track of states that
change sign every time two of them are swapped, and in a random sampling the
positive and negative contributions cancel, leaving noise. That is the most famous
form of the sign problem of Chapter~\ref{ch:history}, and the reason simulating
molecules is one of the problems in Part~\ref{part:problems}.

\subsection*{In two dimensions, any phase}

The argument that left only $\pm1$ hid an assumption: that swapping twice is the
same as doing nothing. Swapping twice means carrying one particle all the way around
the other and back to where it started. In three dimensions, that loop can be lifted
over the other particle and shrunk to a point, so it really is the same as staying
put, and $P^2 = I$. In two dimensions, a loop around a particle cannot leave the
plane. To shrink it, it would have to pass through the particle. Nothing then forces
$P^2 = I$, and the phase of an exchange can be anything.

\begin{figure}[htbp]
\centering
\begin{tikzpicture}[x=1mm, y=1mm]
  % ---------- (a) three dimensions ----------
  \begin{scope}
    \node[font=\sffamily\small] at (20,33) {(a) three dimensions};
    \fill[hfSource] (20,8) circle (1.6);
    \node[font=\small, left=1.2mm] at (20,8) {$B$};
    \draw[hfConcept, line width=1pt] (20,8) ellipse (14 and 4.5);
    \fill[hfConcept] (34,8) circle (1.4);
    \node[font=\small, right=1mm] at (34,8) {$A$};
    \draw[hfConcept, -{Stealth[length=2mm]}, line width=0.8pt] (21,12.5) -- (19.5,12.5);
    % lifted loop
    \draw[hfConcept, dashed, line width=0.8pt] (20,20) ellipse (8 and 2.5);
    \draw[hfConcept, dashed, line width=0.8pt] (20,27) ellipse (2.5 and 0.8);
    \draw[black!40, -{Stealth[length=2mm]}, line width=0.8pt] (37,14) to[bend right=20] (29,24);
    \node[hf/note, anchor=west] at (37,22) {lift the loop\\over $B$: it\\shrinks to\\nothing};
    \node[hf/note] at (20,-5) {going around twice is the same as\\staying put, so $P^2 = I$ and the\\phase is $+1$ or $-1$};
  \end{scope}
  % ---------- (b) two dimensions ----------
  \begin{scope}[xshift=78mm]
    \node[font=\sffamily\small] at (20,33) {(b) two dimensions};
    \fill[black!5] (-6,1) -- (40,1) -- (48,17) -- (2,17) -- cycle;
    \draw[black!30, line width=0.6pt] (-6,1) -- (40,1) -- (48,17) -- (2,17) -- cycle;
    \fill[hfSource] (21,9) circle (1.6);
    \node[font=\small, left=1.2mm] at (21,9) {$B$};
    \draw[hfConcept, line width=1pt] (21,9) ellipse (13 and 4.5);
    \fill[hfConcept] (34,9) circle (1.4);
    \node[font=\small, right=1mm] at (34,9) {$A$};
    \draw[hfConcept, -{Stealth[length=2mm]}, line width=0.8pt] (22,13.5) -- (20.5,13.5);
    \node[hf/note, text=hfAlert] at (21,24) {the loop cannot leave the plane:\\to shrink it, $A$ must cross $B$};
    \node[hf/note] at (21,-5) {going around twice is not the same\\as staying put, and the phase of\\an exchange can be anything};
  \end{scope}
\end{tikzpicture}
\caption{Why only two kinds of particle exist in three dimensions and more exist in
two. Exchanging $A$ and $B$ twice is the same as carrying $A$ once around $B$. (a) In
space, that loop can be lifted over $B$ and shrunk to a point, so a double exchange
does nothing. (b) On a plane, the loop is trapped around $B$, and the double exchange
can leave a phase behind. Particles that use that freedom are anyons.}\label{fig:exchange}
\end{figure}

\begin{history}{1977--1982 \textperiodcentered{} Oslo and Santa Barbara}
Jon Magne Leinaas and Jan Myrheim, in Oslo, found the loophole in 1977 by looking at
the space of all possible paths of identical particles: in one and two dimensions, an
exchange can give any phase. In 1982 Frank Wilczek, at Santa Barbara, built a concrete
example, a charged particle bound to a thin tube of magnetic flux, and named such
particles \emph{anyons}, because the phase can be anything.
\end{history}

A two-dimensional world sounds like a mathematical game, but electrons confined in a
thin layer of semiconductor, in a strong magnetic field, behave as one. Their collective
excitations in the fractional quantum Hall effect, discovered in 1982, are anyons, and in
2020 two groups, in Paris and at Purdue, measured their exchange phase directly.

There is one more step, and it is the one topological quantum computing rests on. If
several anyons share a set of states with the same energy, an exchange does not have to
multiply the state by a number. It can apply a matrix, a unitary on that set of states,
and then the order of the exchanges matters: exchanging $A$ with $B$ and then $B$ with
$C$ is not the same as the other way around. Such particles are \emph{non-abelian}
anyons, after the mathematician Niels Henrik Abel, whose name is attached to operations
whose order does not matter. Braiding non-abelian anyons around each other is a quantum
gate, and Section~\ref{sec:topological} is about the hardware that tries to build them.
